% !TeX root = ../papers/nc-p3.tex
\section{反対階数六・七での片側分岐の排除}
\label{sec:low-rank-branch-exclusion}
\begin{theorem}\label{thm:low-rank-branch-exclusion}
反対切断の階数 $\rho\le7$ の標準正則窓で射影像が高々一次元なら、全ての辺の全射影の
ファイバーは空か一点である。従ってこの階数範囲の全窓は、曲面射影の分類と
系\ref{cor:low-opposite-point-graph-as}の点グラフ判定で被覆される。
\end{theorem}
\begin{proof}[直交補と第四の辺]
残る分岐があれば前節の有理曲線列を取る。反対切断を $\rho$ 次元空間とその双対の完全対合に分解し、
点 $x$ における係数行列の行空間を $W_x$ とする。その次元は $w(x)$ の中央支持次元である。
隣接点対の条件は二つの行空間の直交である。
列の $x_i,x_{i+1}$ ($0\le i\le m$) で一般に次元和が $\rho$ なら、直交補が一致するので
$x_{i+2}$ と $x_{i-1}$ も隣接点対になる。従ってこの四点を循環できる。
$i=0,1$ では $x_{-1}$ または $x_0$ が固定されており、動く $x_1$ を積に含む
四文字閉経路の二つを反復できて矛盾する。$i=m-1,m$ でも逆順に同じである。
内部なら一般点に一意な前点を取って $x_{i+2}=x_{i-2}$ を得る。
動く既約成分の四位置内反復または定点への一致になり、曲線列の構成に反する。
従って全辺で次元和は $\rho$ 未満である。
動く一般点の両端支持は3なので中央支持も3以上である（補題\ref{lem:middle-binary-pencil}）。
$m\ge2$ のため $\rho\le6$ はこれで済む。$\rho=7$ では全ての動く一般点で中央支持は3、
両定点でも中央支持は高々3である。
\end{proof}

\begin{proof}[各順側中心に対応する逆側中心と純関係]
順側中心 $p$ を位置0に置き、$S\subset V_1$ を二次元順支持とする。
最初の直線 $L=\mathbb P(S^\perp)$ の一般点 $q$ で
$w(q)=\sum K_j(q)\otimes A_j$ と書け、三つの $K_j(q)\in V_2$ は独立である。
従って $V_3$ の中央支持は $A_j$ の支持の和で一定であり、三次元である。
それを消す一点を $y\in\mathbb P(V_3^*)$ とする。全 $q\in L$ で $w(q,y)=0$ なので
$w(y)$ の中央支持は $S$ に含まれ、一次元支持は不可能なので $S$ そのものである。
$w(q,y)=0$ を隣接二切断の両方で読むと、$y$ の両端評価と $q$ の対応する両端評価の次元和は
それぞれ6以下である。一般の $q$ の両端階数は3なので $y$ の両端も3以下である。
中央二次元の束の補題により $w(y)$ は $2,2,3$ 型またはその逆、中央鉛筆はKroneckerの $L_2$ 型である。
特にその二つの二文字支持における三次元の縮約像は、二項二次の対称平面と同型である。
この平面は基底を取れば $s_0m_0,s_0m_1+s_1m_0,s_1m_1$ で張られ、
全ての $[\lambda:\mu]$ について $(\lambda s_0+\mu s_1)(\lambda m_0+\mu m_1)$ を含む。

この平面が $R_0$ にあれば $S$ の全文字に終わる純関係がある。
しかし $w(p)$ の末尾からの三次元縮約像は $S$ と高々三次元の中央支持の積にあり、純関係を含む。
実際 $2\times3$ 行列の射影二平面は三次元Segre多様体と $\mathbb P^5$ で交わる。
これは連続純関係の禁止に反する。従って平面は $R_1$ にあり、
$y$ は逆側中心、その前点の直線を $N=\mathbb P(M^\perp)$, $M\subset V_2$ と書く。
$R_1\cap(S\otimes V_2)$ の次元は3なのでこの交叉は上の $S\otimes M$ 内の平面ちょうどである。
よって $w(p)$ も中央支持が $M$ で $L_2$ 型である。
$N$ の一般点の中央支持も三次元でそれを消す反対点は $p$ であるから、逆順にも同じ構成をすると元に戻る。
すなわち位置 $i$ の順側中心と位置 $i-1$ の逆側中心は一対一である。
\end{proof}

\begin{proof}[純関係による列長の制限]
一組 $p,y$ から得た $R_1$ の全点対は「第一点が $L$ 上」または「第二点が $N$ 上」を満たす。
実際、純関係の二つの一次評価の積がパラメータ直線上で恒等零だからである。
位置1に別の順側中心 $p'$ があり、その順支持を $S'\subset V_2$ とすると $S'\cap M=0$ である。
$R_1$ には $M$ の全方向に終わる純関係があり、$R_2$ には $S'$ の全方向に始まる純関係があるからである。
従って $p'$ の後続直線は $N$ と交わらず、$p'$ 自身は $L$ にある必要がある。
これらの結論は位置を回転しても使える。

長さ $m=2$ の曲線列はない。その最後の直線は位置3の逆側中心の前点全体なので、
対をなす位置0の順側中心（元の $p$ とは限らない）から得た上の $N$ の一つである。
$d_2=1$ ゆえ $f_1=2$ で、ある位置1の順側中心 $p'$ からこの $N$ に後続があることになり矛盾する。
一位置に二つの順側中心もない。第一の中心の後続直線は第二のものとは異なるので、
その次の動く像は第二の $N$ に含まれ、次に一定となり $m=2$ になってしまう。

連続する二位置0,1に順側中心があるとする。位置1の中心は $L$ にあるので $f_1>0$、
$d_2=3-f_1\le2$ である。$x_2$ の像は位置1の中心の後続直線にはなれない。
そこでは一般前点が一定であるのに $x_1$ は変動するからである。
位置1の対から来る純関係により $x_3$ はその対の逆側直線 $N'$ 上である。
$m=2$ は排除済みなので $x_3$ は動き、$m=3$, $d_3=1$ である。
$0\le f_2=3d_2-3-d_3$ より $d_2=2$, $f_2=2$ となる。しかし位置2の順側中心からは
$N'$ に後続を持てない。従って順側中心は隣接位置に存在できない。逆側も同じである。
\end{proof}

\begin{proof}[三次曲線の次位置]
そこで位置0の順側中心から出る最初の直線上に新しい順側中心はなく、$f_1=0$, $d_2=3$ である。
$x_2$ は有理正規三次で、$V_2$ を二項三次の完全系とする。位置1,2を独立な二項変数 $x,v$ で読むと、
$L$ に制限した $w$ は $[x,v]A(v)$ と書ける。$[x,v]$ は対角の行列式で、$A(v)$ は
$U=\ker p$, $Y=\ker y$ の二つの三空間に値を持つ二次の行列である。
これは $x=v$ での消滅と、$L$ 上の一般点における位置0,3支持による。
$A$ は定数側の二因子のいずれでも全三空間を張る。

$Y^*$ による $A(v)$ の縮約の三空間を $D\subset S^2(v)\otimes U$ とする。
$\dim S^1 D\le5$ である。実際この積は制限テンソルの反対切断で $x,Y$ を縮約した像であり、
全 $x\in L$ について $W_x\subset W_p^\perp$、$\dim W_p=2$ なので五次元を経由する。
従って二つの独立な $u_0,u_1\in Y^*$ があり、
$A(v)(v_0u_0+v_1u_1)=0$ である。一次関係の二係数が比例すれば $D$ が三次元であることに反する。
一般の $A(v)$ は正確に階数2である。階数1なら、二項多項式環の一意分解で
$A(v)=h(v)a(v)b(v)^t$ と斉次多項式ベクトルの積にできる（例えば原始な一列を取り各列をその倍にする）。
両支持が3なので $a,b$ は各々二次以上となり、総次数2に反する。

$x_3(v)$ は $A(v)$ を位置3から消す。実際、$x_2(v),x_3(v)$ で $w$ は零となり、
$[x,v]$ は独立な $x$ で非零にできる。従って $x_3$ の $Y$ への制限は一般に
$v_0u_0+v_1u_1$ の倍で、$x_3$ の像は平面に入る。
さらに $T\subset V_3$ を $p$ の三次元逆支持とすると、$T$ への射影は
$w(p,x_2(v))$ の $T^*$ 側の核点である。$x_2$ は全空間を張り、$M$ 上の二座標比は非定数なので、
一般にこれは $2\times3$ の $L_2$ 鉛筆の核点である。
具体的に非零の二項パラメータ $(s,t)$ で二行は $(s,t,0),(0,s,t)$ となり、
核は $[t^2:-st:s^2]$ である。従って $x_3$ はこの滑らかな二次曲線上の錐に入る
（錐頂点への定数写像は $m\ge3$ に反する）。整二次錐は平面を含まないので
平面との交叉に入る動く $x_3$ の次数は高々2である。各 $x_j$ が像に双有理であることを使った。

いま $f_2=3d_2-3-d_3=6-d_3>0$ なので位置2にも順側中心がある。
$m=3$ なら終点の位置0に逆側中心があり、その対の順側中心は位置1となって連続禁止に反する。
従って $m\ge4$。位置2の中心の後続直線には $x_3$ の像は一致しない（一般前点は $x_2$ で動く）。
その純関係により $x_4$ は逆側直線に入り、$m=4$ となる。
他方、逆向きにも隣接する中心がないため最終直線の次に逆にたどる曲線の次数は3である。
これは $d_3=3$ を意味し、先の上界に矛盾する。これで階数7でも全ての分岐を排除した。
\end{proof}

\begin{corollary}\label{cor:pure-free-branch-rank-nine}
曲面のない標準正則窓に純関係がないとき、片側分岐があれば $\rho\ge9$ である。
さらに分岐中心 $p$ の $2\times4\times3$ 縮約の二項鉛筆は $L_3$ 型である。
\end{corollary}
\begin{proof}
$p$ の中央支持が3以下なら、三次元の末尾縮約像が $2\times3$ 以下の行列の空間となり、
上のSegreとの交叉で純関係を生じる。従って中央支持は4である。
順支持二次元で作る $4\times3$ 鉛筆は全非零パラメータで階数3でなければならない。
さもなくば末尾の核ベクトルによる縮約が順支持一次元の非零な純関係になる。
全支持性からKronecker分解には零行も正方部分もなく、$L_3$ 型である。
また $L$ の動く一般点で中央支持は3以上である。3なら上の一定反対点 $y$ が得られ、
その $L_2$ 型からいずれかの隣辺に純関係ができる。従って最初の二点の中央支持はともに4であり、
それらの次元和が $\rho$ より小さいことから結論を得る。
\end{proof}
